% !TeX program = xelatex
% 编译：latexmk -xelatex -shell-escape -interaction=nonstopmode -halt-on-error -outdir=build report.tex
% minted 依赖系统 PATH 中的 pygmentize；图片位于同级 figures/ 目录。

\documentclass[UTF8,12pt,a4paper,fontset=fandol]{ctexart}

\usepackage{geometry}
\geometry{left=2.5cm,right=2.5cm,top=2.6cm,bottom=2.6cm}
\usepackage[table]{xcolor}
\usepackage{graphicx}
\usepackage{booktabs}
\usepackage{array}
\usepackage{amsmath}
\usepackage{enumitem}
\usepackage{float}
\usepackage{caption}
\usepackage{fancyhdr}
\usepackage{minted}
\usepackage{underscore}   % 允许正文中直接书写含下划线的标识符
\usepackage[hidelinks]{hyperref}

% ---------- 配色 ----------
\definecolor{PrimaryBlue}{HTML}{1F4E79}
\definecolor{CodeBackground}{HTML}{F2F3F5}   % 代码浅灰衬底
\definecolor{CodeFrame}{HTML}{4F5862}        % 代码深灰边框
\definecolor{TableHeader}{HTML}{DCE6F1}
\definecolor{MutedText}{HTML}{5B6570}

% ---------- 表格 ----------
\newcolumntype{L}[1]{>{\raggedright\arraybackslash}p{#1}}
\newcolumntype{N}[1]{>{\centering\arraybackslash}p{#1}}
\renewcommand{\arraystretch}{1.28}
\setlength{\tabcolsep}{6pt}

% ---------- 代码：minted，浅灰衬底、深灰边框 ----------
\setmonofont{DejaVu Sans Mono}[Scale=0.82]
\setminted{
  frame=single,
  framesep=3mm,               % 衬底与边框向代码四周延展
  rulecolor=\color{CodeFrame},
  bgcolor=CodeBackground,
  fontsize=\small,
  breaklines,
  breakanywhere,
  autogobble,
  tabsize=4,
  showtabs=false,
  showspaces=false
}

% ---------- 标题、页眉与版面 ----------
\ctexset{
  section={format=\Large\bfseries\color{PrimaryBlue}},
  subsection={format=\large\bfseries},
  subsubsection={format=\normalsize\bfseries}
}
\pagestyle{fancy}
\fancyhf{}
\fancyhead[L]{\small 计算机网络实验}
\fancyhead[R]{\small 03 交换机转发}
\fancyfoot[C]{\small\thepage}
\renewcommand{\headrulewidth}{0.3pt}
\setlength{\headheight}{15pt}
\captionsetup{font=small,labelfont=bf}
\setlist{itemsep=0.3em,topsep=0.4em}
\setlength{\parindent}{2em}
\setlength{\parskip}{0.2em}
\setlength{\emergencystretch}{2em}
\raggedbottom

% 插图统一使用该命令：默认占满页宽，两张拓扑图传入 0.75\textwidth
\newcommand{\figref}[3][\textwidth]{%
  \begin{figure}[H]
    \centering
    \includegraphics[width=#1]{#2}
    \caption{#3}
  \end{figure}}

\begin{document}

% ---------- 标题与个人信息（无表格，标题下方横向平铺） ----------
\begin{center}
  {\zihao{2}\bfseries\color{PrimaryBlue} 03 交换机转发实验报告\par}
  \vspace{0.35em}
  {\large 广播网络与交换机转发\par}
  \vspace{0.9em}
  {\small 姓名：吴峥 \hfill 学号：2024K8009909013 \hfill 班级：2406 \hfill 实验日期：2026 年 9 月 17 日\par}
\end{center}
\vspace{0.2em}
\hrule
\vspace{1.1em}

\section*{一、实验目的与完成情况}

本实验在 Mininet 提供的虚拟网络中，用 C 语言在用户态实现二层网络设备，共包含两个部分：

\begin{enumerate}[leftmargin=1.9em]
  \item \textbf{集线器（Hub）}：实现 \texttt{broadcast_packet()}，使节点具备“从任意接口收到帧后向其余所有接口复制转发”的广播能力；验证广播网络的连通性和转发效率，并自建环形拓扑观察广播帧在网络中形成环路、引发广播风暴的现象。
  \item \textbf{交换机（Switch）}：实现 MAC 地址表 \texttt{mac_port_map} 的全部操作（查找、插入、老化清除），并据此实现数据包的定向转发与未知单播广播，最后用 \texttt{iperf} 对比交换机转发与集线器广播的性能差异。
\end{enumerate}

两部分均已通过线上验证：

\figref{../03-hub+switch/hub/online_judge.png}{集线器线上验证结果}

\figref{../03-hub+switch/switch/online_judge.png}{交换机线上验证结果}

\section*{二、实验环境与公共框架}

\subsection*{2.1 实验环境}

{\small\setlength{\parindent}{0pt}
\begin{tabular}{@{}L{0.28\textwidth}L{0.66\textwidth}@{}}
\toprule
\rowcolor{TableHeader}\textbf{组成} & \textbf{本实验使用的配置}\\
\midrule
运行环境 & WSL 2（Windows Subsystem for Linux 2）\\
Linux 发行版 & Ubuntu 24.04.5 LTS\\
Linux 内核 & \texttt{6.18.33.2-microsoft-standard-WSL2}，x86-64\\
编程语言与编译工具 & C、GCC、Make\\
网络虚拟化 & Mininet 2.3.0，\texttt{TCLink} 带带宽链路\\
抓包与测量 & Wireshark / Dumpcap 4.2.2、\texttt{iperf}、\texttt{ping}\\
主机 CPU & Intel(R) Core(TM) i7-14650HX\\
\bottomrule
\end{tabular}}

\subsection*{2.2 用户态转发框架}

Hub 与 Switch 共用一套用户态转发框架：程序为每个名字形如 \texttt{*-eth*} 的接口打开一个 \texttt{AF_PACKET} / \texttt{SOCK_RAW} 套接字，用 \texttt{poll()} 等待任一接口可读，收到数据包后交给 \texttt{handle_packet()} 处理。核心数据结构如下：

\begin{minted}{c}
typedef struct {
    struct list_head iface_list;    // the list of interfaces
    int nifs;                        // number of interfaces
    struct pollfd *fds;                // structure used to poll packets among
                                    // all the interfaces
} ustack_t;

typedef struct {
    struct list_head list;        // list node used to link all interfaces

    int fd;                        // file descriptor for receiving & sending packets
    int index;                    // the index (unique ID) of this interface
    u8    mac[ETH_ALEN];            // mac address of this interface
    char name[16];                // name of this interface
} iface_info_t;
\end{minted}

所有接口对象挂在全局 \texttt{instance->iface_list} 上，\texttt{iface->index} 是内核对接口的唯一编号。发送一个帧由 \texttt{iface_send_packet()} 完成，它把帧的目的 MAC 写入 \texttt{sockaddr_ll}，再通过 \texttt{sendto()} 从指定接口发出：

\begin{minted}{c}
void iface_send_packet(iface_info_t *iface, const char *packet, int len)
{
    struct sockaddr_ll addr;
    memset(&addr, 0, sizeof(struct sockaddr_ll));
    addr.sll_family   = AF_PACKET;
    addr.sll_ifindex  = iface->index;
    addr.sll_halen    = ETH_ALEN;
    addr.sll_protocol = htons(ETH_P_ARP);
    struct ether_header *eh = (struct ether_header *)packet;
    memcpy(addr.sll_addr, eh->ether_dhost, ETH_ALEN);

    if (sendto(iface->fd, packet, len, 0, (const struct sockaddr *)&addr,
                sizeof(struct sockaddr_ll)) < 0) {
        perror("Send raw packet failed");
    }
}
\end{minted}

主循环的关键部分如下，它保证“收到一个包、处理一个包”，并在处理完成后释放 \texttt{packet} 的内存：

\begin{minted}{c}
while (1) {
    int ready = poll(instance->fds, instance->nifs, -1);
    ...
    for (int i = 0; i < instance->nifs; i++) {
        if (instance->fds[i].revents & POLLIN) {
            len = recvfrom(instance->fds[i].fd, buf, ETH_FRAME_LEN, 0, ...);
            ...
            else {
                iface_info_t *iface = fd_to_iface(instance->fds[i].fd);
                ...
                char *packet = malloc(len);
                memcpy(packet, buf, len);
                handle_packet(iface, packet, len);
            }
        }
    }
}
\end{minted}

Hub 与 Switch 的区别集中在 \texttt{handle_packet()}：Hub 对任何帧都广播，Switch 先查 MAC 地址表再决定转发还是广播。

\subsection*{2.3 三节点拓扑}

集线器与交换机两部分都使用 \texttt{three_nodes_bw.py} 描述的三节点带带宽拓扑：

\figref[0.75\textwidth]{figures/topology_three_nodes.png}{三节点带带宽拓扑}

\begin{minted}{text}
h1（10.0.0.1）── 20 Mbit/s ──┐
h2（10.0.0.2）── 10 Mbit/s ──┼── b1 / s1
h3（10.0.0.3）── 10 Mbit/s ──┘
\end{minted}

链路带宽分别如下：\texttt{h1} 的接入链路为 20 Mbit/s，\texttt{h2}、\texttt{h3} 的接入链路各为 10 Mbit/s。

\section*{三、集线器（Hub）的实现与验证}

\subsection*{3.1 \texttt{broadcast_packet()} 的实现}

集线器工作在物理层，它不理解 MAC 地址，只负责把从一个端口收到的信号复制到其余所有端口：

\begin{minted}{c}
// the memory of ``packet'' will be free'd in handle_packet().
void broadcast_packet(iface_info_t *iface, const char *packet, int len)
{
    iface_info_t *entry;
    list_for_each_entry(entry, &(instance->iface_list), list)    // defined in include/list.h
    {
        if (entry->index != iface->index)
        {
            iface_send_packet(entry, packet, len);
        }
    }
}
\end{minted}

实现要点：

\begin{itemize}[leftmargin=1.7em]
  \item \texttt{list_for_each_entry} 是框架提供的链表遍历宏，它通过 \texttt{offsetof} 从内嵌的 \texttt{struct list_head} 反推出 \texttt{iface_info_t} 的地址。
  \item 判断“是否是接收接口”使用 \texttt{entry->index}（内核接口编号），他是内核对接口的唯一编号，并且生命周期覆盖整个运行周期。
\end{itemize}

\subsection*{3.2 连通性验证}

在 Mininet 中为 b1 启动 \texttt{./hub} 后进行连通性测试：

\begin{minted}[fontsize=\footnotesize]{text}
mininet> xterm b1
mininet> pingall
*** Ping: testing ping reachability
b1 -> X X X 
h1 -> *** Error: could not parse ping output: ping: None: Temporary failure in name resolution

X h2 h3 
h2 -> *** Error: could not parse ping output: ping: None: Temporary failure in name resolution

X h1 h3 
h3 -> *** Error: could not parse ping output: ping: None: Temporary failure in name resolution

X h1 h2 
*** Results: 50% dropped (6/12 received)
mininet> h1 ping h2 -c 2
PING 10.0.0.2 (10.0.0.2) 56(84) bytes of data.
64 bytes from 10.0.0.2: icmp_seq=1 ttl=64 time=0.060 ms
64 bytes from 10.0.0.2: icmp_seq=2 ttl=64 time=0.344 ms

--- 10.0.0.2 ping statistics ---
2 packets transmitted, 2 received, 0% packet loss, time 1123ms
rtt min/avg/max/mdev = 0.060/0.202/0.344/0.142 ms
mininet> h1 ping h3 -c 2
...
64 bytes from 10.0.0.3: icmp_seq=1 ttl=64 time=0.132 ms
64 bytes from 10.0.0.3: icmp_seq=2 ttl=64 time=0.093 ms
rtt min/avg/max/mdev = 0.093/0.112/0.132/0.019 ms
mininet> h3 ping h2 -c 2
...
64 bytes from 10.0.0.2: icmp_seq=1 ttl=64 time=0.054 ms
64 bytes from 10.0.0.2: icmp_seq=2 ttl=64 time=0.070 ms
rtt min/avg/max/mdev = 0.054/0.062/0.070/0.008 ms
\end{minted}

\texttt{pingall} 只有指向 \texttt{b1} 的几处失败。原因是 \texttt{b1} 作为二层设备没有配置 IP 地址，Mininet 的 \texttt{pingall} 仍然尝试 ping 它的 IP，于是得到 \texttt{Temporary failure in name resolution}，这是预期行为。

三个端节点之间的往返时延汇总如下：

\begin{center}
{\small
\begin{tabular}{@{}lrrrr@{}}
\toprule
\rowcolor{TableHeader}\textbf{方向} & \textbf{min RTT} & \textbf{avg RTT} & \textbf{max RTT} & \textbf{丢包}\\
\midrule
h1 → h2 & 0.060 ms & 0.202 ms & 0.344 ms & 0\%\\
h1 → h3 & 0.093 ms & 0.112 ms & 0.132 ms & 0\%\\
h3 → h2 & 0.054 ms & 0.062 ms & 0.070 ms & 0\%\\
\bottomrule
\end{tabular}}
\end{center}

时延都在亚毫秒级，说明虚拟链路上的广播转发路径通畅。

\subsection*{3.3 广播网络的效率测量}

在 \texttt{three_nodes_bw.py} 上使用 \texttt{iperf} 测量两种场景：

\begin{itemize}[leftmargin=1.7em]
  \item \textbf{场景 A（H1 为发送端）}：\texttt{H1} 同时作为 client 向 \texttt{H2}、\texttt{H3} 两个 server 发送数据；
  \item \textbf{场景 B（H1 为接收端）}：\texttt{H1} 作为 server，\texttt{H2}、\texttt{H3} 两个 client 同时向 \texttt{H1} 发送数据。
\end{itemize}

场景 A 与场景 B 的原始日志如下（每段测试 30 秒，TCP 端口 5001，窗口 85.3 KByte）：

\begin{minted}[fontsize=\footnotesize]{bash}
# 场景 A
iperf -c 10.0.0.2 -t 30 & iperf -c 10.0.0.3 -t 30
[1] 61416
------------------------------------------------------------
Client connecting to 10.0.0.3, TCP port 5001
Client connecting to 10.0.0.2, TCP port 5001
TCP window size: 85.3 KByte (default)TCP window size: 85.3 KByte (default)
------------------------------------------------------------
[  1] local 10.0.0.1 port 52310 connected with 10.0.0.3 port 5001 (icwnd/mss/irtt=14/1448/253)
[  1] local 10.0.0.1 port 42300 connected with 10.0.0.2 port 5001 (icwnd/mss/irtt=14/1448/300)
[ ID] Interval       Transfer     Bandwidth
[  1] 0.0000-30.4909 sec  21.3 MBytes  5.85 Mbits/sec
[  1] 0.0000-30.8921 sec  11.0 MBytes  2.99 Mbits/sec
\end{minted}

\begin{minted}[fontsize=\footnotesize]{bash}
# 场景 B
iperf -s
------------------------------------------------------------
Server listening on TCP port 5001
TCP window size: 85.3 KByte (default)
------------------------------------------------------------
[  1] local 10.0.0.1 port 5001 connected with 10.0.0.3 port 35392 (icwnd/mss/irtt=14/1448/140)
[  2] local 10.0.0.1 port 5001 connected with 10.0.0.2 port 44322 (icwnd/mss/irtt=14/1448/21686)
[ ID] Interval       Transfer     Bandwidth
[  1] 0.0000-30.8588 sec  30.4 MBytes  8.26 Mbits/sec
[  2] 0.0000-30.6453 sec  30.6 MBytes  8.38 Mbits/sec
\end{minted}

整理成表格：

\begin{center}
{\small
\begin{tabular}{@{}llrrrr@{}}
\toprule
\rowcolor{TableHeader}\textbf{场景} & \textbf{数据流} & \textbf{传输量} & \textbf{时长} & \textbf{吞吐} & \textbf{合计}\\
\midrule
A & 流 1 & 21.3 MB & 30.49 s & 5.85 Mbit/s & \textbf{8.84 Mbit/s}\\
  & 流 2 & 11.0 MB & 30.89 s & 2.99 Mbit/s & \\
B & H2 → H1 & 30.6 MB & 30.65 s & 8.38 Mbit/s & \textbf{16.64 Mbit/s}\\
  & H3 → H1 & 30.4 MB & 30.86 s & 8.26 Mbit/s & \\
\bottomrule
\end{tabular}}
\end{center}

可以看到两个明显现象：

\begin{enumerate}[leftmargin=1.9em]
  \item \textbf{场景 A 的合计吞吐（8.84 Mbit/s）只有场景 B（16.64 Mbit/s）的约 53\%}，广播网络在“一个节点同时发给两个节点”时效率明显更差。
  \item \textbf{同为发送端，两条流的带宽分配严重不均}：一条得到 5.85 Mbit/s，另一条只有 2.99 Mbit/s，接近 2:1。
\end{enumerate}

原因在第五章结合交换机结果一起分析。

\subsection*{3.4 环形拓扑下的广播环路}

\subsubsection*{（1）环形拓扑的构造}

\texttt{ring_topo.py} 在实验一的框架上增加了第三个 Hub \texttt{b3}，让 \texttt{b1} 与 \texttt{b2} 之间同时存在两条路径：一条是直连的 \texttt{b1—b2}，另一条是绕行 \texttt{b1—b3—b2}。于是三台 Hub 与链路一起构成了一个物理环路：

\figref[0.75\textwidth]{figures/topology_ring.png}{环形拓扑：三台 Hub 之间的环路}

\begin{minted}{python}
#            h1 --- b1 --- b2 --- h2
#                    \     /
#                     \   /
#                      b3
#
# b1 和 b2 之间有两条路：直连的 b1--b2，以及绕行的 b1--b3--b2。
# 三台 hub 一起跑起来之后，一个广播帧会沿着 b1->b2->b3->b1 永远转圈，
# 而且每经过一个 hub 就被复制一份。
class RingTopo(Topo):
    def build(self):
        h1 = self.addHost('h1')
        h2 = self.addHost('h2')
        b1 = self.addHost('b1')
        b2 = self.addHost('b2')
        b3 = self.addHost('b3')

        self.addLink(h1, b1, bw=20)     # h1-eth0  <-> b1-eth0
        self.addLink(b1, b2, bw=10)     # b1-eth1  <-> b2-eth0
        self.addLink(b2, h2, bw=20)     # b2-eth1  <-> h2-eth0
        self.addLink(b1, b3, bw=10)     # b1-eth2  <-> b3-eth0   ┐
        self.addLink(b3, b2, bw=10)     # b3-eth1  <-> b2-eth2   ┘ 环路: 绕行那条
\end{minted}

接口分布如下：

\begin{minted}{text}
b1: eth0->h1   eth1->b2   eth2->b3
b2: eth0->b1   eth1->h2   eth2->b3
b3: eth0->b1   eth1->b2
\end{minted}

\subsubsection*{（2）环路抓包与统计}

在环形拓扑上执行一次 \texttt{ping -c 1}，并在 \texttt{h1-eth0} 上抓包，得到结果整理如下：

\begin{center}
{\small
\begin{tabular}{@{}ll@{}}
\toprule
\rowcolor{TableHeader}\textbf{项目} & \textbf{数值}\\
\midrule
抓包接口 & \texttt{h1-eth0}\\
抓包工具 & Dumpcap (Wireshark) 4.2.2\\
数据包数 & 206,188\\
抓包时长 & 6.245751322 s\\
总数据量 & 约 15 MB\\
平均包长 & 74.91 Byte\\
平均包速率 & 33,012 packet/s（33 kpps）\\
平均码率 & 19.78 Mbit/s\\
\bottomrule
\end{tabular}}
\end{center}

\begin{center}
{\small
\begin{tabular}{@{}lr@{}}
\toprule
\rowcolor{TableHeader}\textbf{指标} & \textbf{数值}\\
\midrule
帧总数 & 206,188\\
\textbf{不重复帧数} & \textbf{487}\\
\textbf{放大倍数（总帧数 / 不重复帧数）} & \textbf{423.4×}\\
ARP 帧 & 85,014（41.2\%）\\
ICMP Echo Reply 帧 & 120,691（58.5\%）\\
ICMP Echo Request 帧 & 483（0.2\%）\\
不重复的 Echo Request（按 id + seq） & \textbf{1}\\
ARP 请求 / ARP 应答 & 350 / 84,664\\
IP 方向统计 & h1→h2：483；h2→h1：120,691\\
\bottomrule
\end{tabular}}
\end{center}

\figref{figures/ring_frame_composition.png}{广播风暴中的帧类型构成}

\figref{figures/ring_broadcast_storm_timeseries.png}{环形拓扑下广播风暴的速率与吞吐随时间变化}

\figref{figures/ring_frame_sizes.png}{广播风暴中的帧长分布}

\subsubsection*{（3）现象解释}

一个\textbf{只发了一次}的 ICMP Echo Request，在 6.2 秒的抓包里出现了 483 次，并且 487 个不重复帧被复制成了 206,188 个帧，放大 423 倍，速率稳定在约 33 kpps、19.8 Mbit/s，直到抓包结束都没有衰减。这正是“广播环路”造成“广播风暴”的典型表现，其形成机制是：

\begin{enumerate}[leftmargin=1.9em]
  \item \textbf{Hub 没有学习能力}。\texttt{broadcast_packet()} 对任何帧都从其余所有接口发出，因此一个帧进入环路后不会被“记住方向”而被丢弃。
  \item \textbf{环路上存在两条路径}。\texttt{b1—b2} 直连与 \texttt{b1—b3—b2} 绕行让同一个帧可以沿两个方向同时前进；每经过一个 Hub 就被复制一份，于是帧的数量按回路不断翻倍。
  \item \textbf{二层转发没有 TTL 机制}。以太网帧头没有类似 IP TTL 的跳数限制，交换机/Hub 也不会递减任何计数，因此帧不会“过期”。
  \item \textbf{没有生成树协议（STP）}。真实交换机会通过 STP 在逻辑上阻塞冗余链路、消除环路，而本实验的 Hub 只做无条件复制。
\end{enumerate}

抓包还给出了两个特征：

\begin{itemize}[leftmargin=1.7em]
  \item Echo Request 只有 \textbf{1 个不重复签名}，却出现 483 次，说明这 483 个请求其实是\textbf{同一份帧在环路中的副本}；\texttt{h2} 每收到一份副本就回一个应答，于是产生了 120,691 个 Echo Reply。一次 ping 变成了双方共同参与的指数级放大。
  \item IP 方向统计出现严重不对称：\texttt{h1→h2} 只有 483 个请求，而 \texttt{h2→h1} 有 120,691 个应答。应答本身也在环路中不断复制，并且因为目的地址是 \texttt{h1}，其中相当一部分副本最终被送到 \texttt{h1} 的链路上，\texttt{h1} 因此反复收到同一次 ping 的重复应答。
\end{itemize}

从时间序列看，包速率与吞吐在整个抓包窗口内几乎保持在 20 Mbit/s 的平台上：\texttt{h1-eth0} 的链路容量恰好是 20 Mbit/s，说明\textbf{广播风暴已经把 \texttt{h1} 的接入链路完全占满}，可用于正常通信的带宽接近于零。值得注意的是，抓包中的帧长只有 98 Byte（ICMP，121,174 帧）和 42 Byte（ARP，85,014 帧）两种，风暴中几乎没有大帧，却仍然能把链路塞满，可见环路对带宽的消耗有多么迅速。

\section*{四、交换机（Switch）的实现}

交换机的核心是“自学习 + 定向转发”：维护一张“MAC 地址 → 接口”的映射表，收到帧后先按目的 MAC 查表，查到就从对应接口转发，查不到才广播；同时把帧的源 MAC 与接收接口登记到表中，并周期性地清除长期不用的表项。

\subsection*{4.1 数据结构与哈希}

\begin{minted}{c}
#define MAC_PORT_TIMEOUT 30

struct mac_port_entry {
    struct list_head list;
    uint8_t mac[ETH_ALEN];
    iface_info_t *iface;
    time_t visited;
};

typedef struct {
    struct list_head hash_table[HASH_8BITS];
    pthread_mutex_t lock;
    pthread_t thread;
} mac_port_map_t;
\end{minted}

MAC 地址表被组织成 256 个桶的哈希表（\texttt{HASH_8BITS = 256}），每个桶是一条链表，表项记录 MAC、出接口和最近一次访问时间 \texttt{visited}。哈希函数直接复用框架提供的 \texttt{hash8()}（逐字节异或）：

\begin{minted}{c}
static inline u8 hash8(char *buf, int len)
{
    u8 result = 0;
    for (int i = 0; i < len; i++)
        result ^= buf[i];

    return result;
}
\end{minted}

在此基础上定义两个宏，使后续代码更易读：

\begin{minted}{c}
// hash: u8 mac[ETH_ALEN] -> HASH_8BITS
#define mac_to_addr_8(mac) (hash8((char *)(mac), ETH_ALEN))

// check if two mac addrs are equal
#define mac_equal(mac1, mac2) (memcmp((mac1), (mac2), ETH_ALEN) == 0)
\end{minted}

\subsection*{4.2 \texttt{lookup_port()}：查找出接口}

\begin{minted}{c}
iface_info_t *lookup_port(u8 mac[ETH_ALEN])
{
    u8 hash_mac = mac_to_addr_8(mac);
    iface_info_t* iface = NULL;

    pthread_mutex_lock(& mac_port_map.lock);
    mac_port_entry_t *entry;
    list_for_each_entry(entry, & mac_port_map.hash_table[hash_mac], list)
    {
        if (mac_equal(mac, entry->mac)) 
        {
            iface = entry->iface;
            entry->visited = time(NULL);
            break;
        }
    }
    pthread_mutex_unlock(&mac_port_map.lock);

    return iface;
}
\end{minted}

函数先算出桶号，只遍历该桶；命中后返回出接口，并\textbf{顺手把 \texttt{visited} 更新为当前时间}。这一点很重要：只要某个 MAC 还在参与通信，它的表项就会不断被刷新，不会被老化线程误删。查不到时返回 \texttt{NULL}，由上层决定改用广播。

\subsection*{4.3 \texttt{insert_mac_port()}：学习源 MAC}

\begin{minted}{c}
void insert_mac_port(u8 mac[ETH_ALEN], iface_info_t *iface)
{
    u8 hash_mac = mac_to_addr_8(mac);

    pthread_mutex_lock(& mac_port_map.lock);
    mac_port_entry_t *entry;
    list_for_each_entry(entry, & mac_port_map.hash_table[hash_mac], list)
    {
        if (mac_equal(mac, entry->mac)) 
        {
            list_delete_entry(& entry->list);
            entry->iface = iface;
            entry->visited = time(NULL);
            list_add_head(& entry->list, & mac_port_map.hash_table[hash_mac]);
            pthread_mutex_unlock(&mac_port_map.lock);
            return;
        }
    }
    mac_port_entry_t *new = malloc (sizeof (mac_port_entry_t));
    if (!new) { pthread_mutex_unlock(& mac_port_map.lock); return; }

    memcpy(new->mac, mac, ETH_ALEN);
    new->iface = iface;
    new->visited = time(NULL);
    list_add_head(& new->list, & mac_port_map.hash_table[hash_mac]);
    pthread_mutex_unlock(& mac_port_map.lock);
}
\end{minted}

处理逻辑分两种情况：

\begin{itemize}[leftmargin=1.7em]
  \item \textbf{表项已存在}：无论是否有接口变化（因为更新接口开销实际并不大，不会形成瓶颈），先把它从原链表位置摘下，更新出接口和时间后重新插入到桶头。这里在 \texttt{list_for_each_entry} 中删除并立即 \texttt{return}，不会继续遍历被修改的链表，因此无需使用 safe 版本。
  \item \textbf{表项不存在}：\texttt{malloc} 一个新表项，填入 MAC、出接口和时间，插入桶头。\texttt{malloc} 失败时只解锁并返回，不影响后续转发——学习失败最多让下一帧继续广播，不会造成崩溃。
\end{itemize}

新表项插入桶头（\texttt{list_add_head}）而不是桶尾，是因为活跃表项会被反复刷新并移到桶头，长期不用的表项自然沉到桶尾，链表局部性更好。

\subsection*{4.4 \texttt{sweep_aged_mac_port_entry()}：老化清除}

\begin{minted}{c}
int sweep_aged_mac_port_entry()
{
    int count = 0;
    time_t now = time(NULL);

    pthread_mutex_lock(&mac_port_map.lock);
    mac_port_entry_t *entry, *q;
    
    for (int i = 0; i < HASH_8BITS; i++) 
    {
        list_for_each_entry_safe(entry, q, &mac_port_map.hash_table[i], list) 
        {
            if (difftime(now, entry->visited) >= MAC_PORT_TIMEOUT)
            {
                list_delete_entry(&entry->list);
                free(entry);
                count ++;
            }
        }
    }

    pthread_mutex_unlock(&mac_port_map.lock);
    return count;
}
\end{minted}

由于遍历过程中会删除并释放当前节点，必须使用 \texttt{list_for_each_entry_safe}。函数返回本次删除的表项数量，供调用者记录日志。

老化线程每秒执行一次清除：

\begin{minted}{c}
void *sweeping_mac_port_thread(void *nil)
{
    while (1) {
        sleep(1);
        int n = sweep_aged_mac_port_entry();

        if (n > 0)
            log(DEBUG, "%d aged entries in mac_port table are removed.", n);
    }

    return NULL;
}
\end{minted}

定期老化可以解决两类问题：一是主机下线后表项永久占用内存；二是主机更换接口时，若旧表项还在，会短时间把帧错误转发到旧端口。30 秒超时是“及时性”与“稳定性”之间的折中。

\subsection*{4.5 \texttt{handle_packet()}：转发与广播}

\begin{minted}{c}
void handle_packet(iface_info_t *iface, char *packet, int len)
{
    struct ether_header *eh = (struct ether_header *)packet;

    iface_info_t *out = lookup_port(eh->ether_dhost);
    if (out)
        iface_send_packet(out, packet, len);        // send to dst
    else
        broadcast_packet(iface, packet, len);       // broadcast

    insert_mac_port(eh->ether_shost, iface);
    
    free(packet);
}
\end{minted}

这是交换机的数据平面，三步顺序与实验注释完全对应：

\begin{enumerate}[leftmargin=1.9em]
  \item \textbf{查表转发}：以目的 MAC 查表。查到则直接单播到对应接口，不再打扰其他端口；查不到（未知单播）或目的地址是广播/组播地址时，调用 \texttt{broadcast_packet()} 从其余接口泛洪。
  \item \textbf{学习源地址}：把源 MAC 与接收接口写入表中。注意必须在转发之后学习，避免把本帧的源地址立刻用于本帧的回环判定——本框架中同一帧不会被发回接收接口，顺序其实不影响正确性，但“先转发、后学习”更符合标准交换机的处理流程。
  \item \textbf{释放内存}：框架在收到帧时 \texttt{malloc}，处理完毕后由 \texttt{handle_packet()} 统一 \texttt{free()}，\texttt{broadcast_packet()} 只借用缓冲区。
\end{enumerate}

交换机的 \texttt{broadcast_packet()} 与 Hub 的完全一致，只是当 MAC 表命中时它不会再被调用，因此泛洪只发生在网络启动初期、表项老化之后或真正的广播帧上：

\begin{minted}{c}
void broadcast_packet(iface_info_t *iface, const char *packet, int len)
{
    iface_info_t *entry;
    list_for_each_entry(entry, &(instance->iface_list), list)
    {
        if (entry->index != iface->index)
        {
            iface_send_packet(entry, packet, len);
        }
    }
}
\end{minted}

\subsection*{4.6 连通性与性能验证}

交换机版本的连通性测试结果：

\begin{minted}{text}
mininet> pingall
*** Ping: testing ping reachability
h1 -> h2 h3 X 
h2 -> h1 h3 X 
h3 -> h1 h2 X 
s1 -> X X X 
*** Results: 50% dropped (6/12 received)
\end{minted}

与 Hub 相同，6 个端节点互 ping 全部成功，失败的 6 个方向都是指向没有 IP 的 \texttt{s1}。

性能测量使用与 Hub 完全相同的拓扑和命令，得到：

\begin{minted}[fontsize=\footnotesize]{bash}
# 场景 A：h1 中执行 iperf -c 10.0.0.2 -t 30 & iperf -c 10.0.0.3 -t 30
[  1] local 10.0.0.1 port 51462 connected with 10.0.0.2 port 5001 (icwnd/mss/irtt=14/1448/166)
[  1] local 10.0.0.1 port 55418 connected with 10.0.0.3 port 5001 (icwnd/mss/irtt=14/1448/118)
[ ID] Interval       Transfer     Bandwidth
[  1] 0.0000-31.0876 sec  34.8 MBytes  9.38 Mbits/sec
[  1] 0.0000-31.2286 sec  35.0 MBytes  9.40 Mbits/sec
\end{minted}

\begin{minted}[fontsize=\footnotesize]{bash}
# 场景 B：h1 中执行 iperf -s，h2、h3 中分别执行 iperf -c 10.0.0.1 -t 30
[  1] local 10.0.0.1 port 5001 connected with 10.0.0.2 port 42756 (icwnd/mss/irtt=14/1448/28)
[  2] local 10.0.0.1 port 5001 connected with 10.0.0.3 port 39966 (icwnd/mss/irtt=14/1448/49)
[ ID] Interval       Transfer     Bandwidth
[  1] 0.0000-31.0278 sec  34.4 MBytes  9.29 Mbits/sec
[  2] 0.0000-30.9100 sec  34.3 MBytes  9.30 Mbits/sec
\end{minted}

\begin{center}
{\small
\begin{tabular}{@{}llrrrr@{}}
\toprule
\rowcolor{TableHeader}\textbf{场景} & \textbf{数据流} & \textbf{传输量} & \textbf{时长} & \textbf{吞吐} & \textbf{合计}\\
\midrule
A & 流 1 & 34.8 MB & 31.09 s & 9.38 Mbit/s & \textbf{18.78 Mbit/s}\\
A & 流 2 & 35.0 MB & 31.23 s & 9.40 Mbit/s & \\
B & H2 → H1 & 34.4 MB & 31.03 s & 9.29 Mbit/s & \textbf{18.59 Mbit/s}\\
B & H3 → H1 & 34.3 MB & 30.91 s & 9.30 Mbit/s & \\
\bottomrule
\end{tabular}}
\end{center}

交换机下两种场景的合计吞吐几乎相同（18.78 与 18.59），且每条流都接近各自 10 Mbit/s 接入链路的理论上限，两条流的分配也非常均匀。

\section*{五、性能对比与分析}

\subsection*{5.1 结果汇总}

\figref{figures/perf_hub_vs_switch.png}{集线器与交换机在两种场景下的分场景吞吐}

\figref{figures/perf_aggregate_speedup.png}{合计吞吐与提速倍数}

\begin{center}
{\small
\begin{tabular}{@{}llrrr@{}}
\toprule
\rowcolor{TableHeader}\textbf{场景} & \textbf{数据流} & \textbf{Hub} & \textbf{Switch} & \textbf{Switch / Hub}\\
\midrule
A & 流 1 & 5.85 & 9.38 & 1.60×\\
A & 流 2 & 2.99 & 9.40 & 3.14×\\
A & \textbf{合计} & \textbf{8.84} & \textbf{18.78} & \textbf{2.12×}\\
B & H2 → H1 & 8.38 & 9.29 & 1.11×\\
B & H3 → H1 & 8.26 & 9.30 & 1.13×\\
B & \textbf{合计} & \textbf{16.64} & \textbf{18.59} & \textbf{1.12×}\\
\bottomrule
\end{tabular}}
\end{center}

换算成对 \texttt{h1} 接入链路（20 Mbit/s）的利用率：

\begin{center}
{\small
\begin{tabular}{@{}lrr@{}}
\toprule
\rowcolor{TableHeader}\textbf{指标} & \textbf{Hub} & \textbf{Switch}\\
\midrule
场景 A 合计吞吐 / H1 链路容量 & 44.2\% & 93.9\%\\
场景 B 合计吞吐 / H1 链路容量 & 83.2\% & 93.0\%\\
场景 B 与场景 A 的合计之比 & 1.88 & 0.99\\
\bottomrule
\end{tabular}}
\end{center}

\subsection*{5.2 场景 A 差异分析：泛洪让下行链路被两条流共享}

场景 A 中 \texttt{H1} 同时向 \texttt{H2}、\texttt{H3} 发送数据。对 \textbf{Hub} 而言，\texttt{broadcast_packet()} 将每个帧发送给所有非传入端口的端口，因此：

\begin{itemize}[leftmargin=1.7em]
  \item \texttt{H1 → H2} 的数据帧不仅出现在 \texttt{H2} 的 10 Mbit/s 链路上，也被复制到 \texttt{H3} 的 10 Mbit/s 链路上；
  \item 同理，\texttt{H1 → H3} 的数据帧也出现在 \texttt{H2} 的链路上。
\end{itemize}

也就是说，\textbf{每一条 10 Mbit/s 的下行接入链路都要同时承载两条流的数据}，其带宽上限直接决定了“两条流之和”的上限。两条流合计 8.84 Mbit/s，已接近 10 Mbit/s 链路在扣除 TCP 首部与 ACK 之后的有效容量，这正是 Hub 场景 A 只能跑到约 8.8 Mbit/s 的原因。同时，因为两条流共享同一条受限链路，TCP 的竞争使带宽分配失衡，出现了 5.85 : 2.99 的不公平划分。

对 \textbf{Switch} 而言，转发大部分时间 \texttt{lookup_port()} 命中，数据帧只会送到目的主机所在的接口：\texttt{H1 → H2} 只走 \texttt{H2} 的链路，\texttt{H1 → H3} 只走 \texttt{H3} 的链路。因此每条 10 Mbit/s 链路只承载一条流，各自可以跑满约 9.4 Mbit/s；而两条流之和 18.78 Mbit/s 已经接近 \texttt{H1} 的 20 Mbit/s 链路容量（93.9\%）。场景 A 的提升倍数达到 \textbf{2.12×}，本质是“把被泛洪浪费掉的另一条下行链路还给了用户”。

\subsection*{5.3 场景 B 差异分析：泛洪副本落在空闲方向上}

场景 B 中 \texttt{H2}、\texttt{H3} 同时向 \texttt{H1} 发送数据，此时：

\begin{itemize}[leftmargin=1.7em]
  \item 每个发送方的\textbf{上行方向}只承载自己的数据，Hub 的复制副本落到的是另一台主机的\textbf{下行方向}，与发送方的上行方向互不争用；
  \item 接收方 \texttt{H1} 的链路容量为 20 Mbit/s，两条流合计 16.64 Mbit/s 仍在容量之内。
\end{itemize}

因此 Hub 在场景 B 下没有被“下游共享”直接卡死，实测 16.64 Mbit/s，与 Switch 的 18.59 Mbit/s 只差约 11\%。这 11\% 的差距来自泛洪本身的代价：Hub 需要为每个帧额外复制、发送一份到无关主机，占用了一份额外的链路与转发开销（交换机只需发送一次），同时无关主机的协议栈也在处理本不属于自己的数据。Switch 侧因为每个帧只发送一次，两条流都能各自跑满约 9.3 Mbit/s。

\subsection*{5.4 广播网络的对称性差异}

把两种场景放在一起看，可以得到一个更直观的结论：

\begin{itemize}[leftmargin=1.7em]
  \item \textbf{交换机是对称的}：场景 A 18.78 Mbit/s，场景 B 18.59 Mbit/s，两者之比 0.99。这是因为转发路径与“谁是发送方”无关，每条流的路径都是“源接入链路 → s1 → 目的接入链路”，与拓扑完全对称。
  \item \textbf{集线器是明显不对称的}：场景 B 是场景 A 的 1.88 倍。因为 Hub 的性能取决于“广播副本落在哪里”：当副本落在已经被占用的下行链路时（场景 A），性能损失巨大；当副本恰好落在空闲方向时（场景 B），损失就小得多。这也指出了 Hub 的吞吐与流量方向紧密相关。
\end{itemize}

\section*{六、实验总结}

本实验完成了广播网络与交换式网络两种用户态二层设备的实现：

\begin{enumerate}[leftmargin=1.9em]
  \item \textbf{Hub}：\texttt{broadcast_packet()} 通过遍历接口链表、排除接收接口，实现了从所有其他端口复制转发。环形拓扑实验定量地展示了缺少学习与环路控制的后果。
  \item \textbf{Switch}：实现了基于 256 桶哈希表的 MAC 地址表，\texttt{lookup_port()} 查找并按访问时间刷新、\texttt{insert_mac_port()} 学习并更新源地址、\texttt{sweep_aged_mac_port_entry()} 每 30 秒清除老化表项，\texttt{handle_packet()} 完成“命中单播、未命中泛洪、随后学习”的完整转发流程；数据平面与老化线程之间用互斥锁保证线程安全，销毁时先停线程再释放资源。
  \item \textbf{性能对比}：在完全相同的拓扑与命令下，交换机的合计吞吐为 18.78 / 18.59 Mbit/s（场景 A / B），接近 \texttt{h1} 链路 20 Mbit/s 容量的 94\%；集线器只有 8.84 / 16.64 Mbit/s，分别为 44\% 和 83\%。\textbf{场景 A 提升 2.12 倍，场景 B 提升 1.12 倍。}
\end{enumerate}

实验指出集线器对每个帧无条件复制到所有其他端口，导致多个流共享同一条下行链路并浪费带宽；交换机通过 MAC 地址学习把单播帧只送到目的端口，让每条流独享自己的链路。实验也说明，二层网络若要引入冗余链路提高可靠性，就必须用 STP 之类的机制消除环路，否则广播帧会立刻演化为风暴。

\end{document}
